% ============================================================
% sec/02trabajos_relacionados.tex — VERSIÓN INTEGRADA, TRES EJES
%
% Eje A: Duarte · Eje B: Chaves · Eje C: Hernández
%
% SECCIÓN CENTRAL DE LA ETAPA 1.
% Tres ejes + tabla comparativa obligatoria + brechas en
% apartado propio de cierre.
% ============================================================
\section{Trabajos Relacionados}
\label{sec:relacionados}

Los candidatos se recuperaron de Scopus, IEEE Xplore, la ACM Digital
Library, Springer, los \textit{proceedings} de AAAI y AIIDE, y los registros
de los editores de las revistas correspondientes. Las cadenas de búsqueda
combinaron \code{win prediction}, \code{game state}, \code{collectible card
game} y \code{esport replay dataset} para el dominio, y \code{data leakage},
\code{grouped cross-validation}, \code{repeated measures},
\code{temporal dataset shift} y \code{tabular deep learning} para la
metodología.

La inclusión exigió venue arbitrado de revista o conferencia, un DOI que
resolviera, y contenido que respaldara una decisión efectivamente tomada en
este trabajo. Se excluyeron los estudios exclusivamente previos a la partida
cuando no aportaban a la representación de estados ni a la generalización.
La mayoría de los trabajos retenidos se publicó en 2021 o después; los más
antiguos se conservan de forma deliberada por ser referencias fundacionales
del campo o el único antecedente publicado de la tarea.

\pc{Completar con el número total de trabajos cribados y retenidos por el
equipo, consolidando los registros de selección de los tres ejes.}

\razon{Este párrafo convierte la revisión en un procedimiento auditable en
lugar de una lista de lecturas. Hodge \textit{et al.} organizan su revisión
declarando los temas que la estructuran \cite{hodge_winprediction_2021}, y
Grinsztajn \textit{et al.} enuncian criterios explícitos de inclusión de
conjuntos de datos \cite{grinsztajn_trees_2022}; ambas prácticas son el
precedente de exigir transparencia en el procedimiento.}

% EJE A

\subsection{Eje A: el problema y su dominio}
\label{subsec:ejea}

\subsubsection{Predicción del desenlace a partir de un único estado}

La formulación dominante trata la predicción del desenlace como
clasificación binaria supervisada sobre una instantánea de la posición, con
la etiqueta tomada del resultado registrado de la partida completa. El
AAIA'17 Data Mining Challenge lo enuncia exactamente en esos términos para
\emph{Hearthstone}: los participantes recibían la descripción de un estado de
juego y entregaban la verosimilitud de que el jugador uno venciera, puntuada
mediante el área bajo la curva ROC \cite{janusz_hearthstone_2017}. El
descriptor SC2EGSet aplica el mismo encuadre a \emph{StarCraft~II},
derivando indicadores económicos de las repeticiones y clasificando el
desenlace a partir de ellos \cite{bialecki_sc2egset_2023}. Ninguna de las dos
formulaciones requiere un simulador en el momento de la predicción, y ninguna
trata la trayectoria como un problema de control. Esta es la formulación que
se adopta aquí.

La revisión sistemática de Kamal \textit{et al.} sitúa estos trabajos en un
campo más amplio: incluye treinta y cinco estudios sobre juegos de arena
multijugador y encuentra un predominio del aprendizaje supervisado, con uso
frecuente de estadísticas de jugador, mecánicas del juego y eventos
intra-partida \cite{kamal_moba_2025}.

Lo que estos trabajos establecen es que la formulación es viable y que la
señal es sustancial. Lo que no establecen es que el nivel alcanzable se
transfiera entre juegos. Los puntos de operación reportados difieren más
entre sí que la diferencia entre métodos dentro de cualquiera de los
estudios, y no están medidos en la misma escala: Janusz \textit{et al.}
reportan área bajo la curva sobre estados extraídos de partidas entre agentes
que eligen sus acciones al azar, mientras que Bia\l ecki \textit{et al.}
reportan acierto sobre características promediadas a lo largo de toda la
partida.

\subsubsection{Capacidad predictiva a lo largo de la partida}

Un segundo tema, presente en varios trabajos pero desarrollado por ninguno,
es que la capacidad predictiva no es un escalar.

Janusz \textit{et al.} grafican el área bajo la curva sobre turnos
consecutivos y muestran que crece a lo largo de la partida
\cite{janusz_hearthstone_2017}. Bia\l ecki \textit{et al.} grafican el
acierto sobre el tiempo de partida y observan que sus tres clasificadores
alcanzan una asíntota alrededor de los 550 segundos, lo que interpretan como
el punto en que los indicadores económicos dejan de aportar información
adicional \cite{bialecki_sc2egset_2023}. Tian \textit{et al.} encuentran que
la utilidad de las variables cambia con la etapa: la alineación domina al
inicio, el oro gana relevancia en la etapa media, y las bajas y las torres se
vuelven más informativas posteriormente; reportan una exactitud del
81.8\,\% en el nodo temporal de siete minutos según su tabla detallada, si
bien el resumen del artículo menciona 84.7\,\% \cite{tian_hero_2022}.

\razon{La discrepancia entre el valor de la tabla y el del resumen debe
consignarse. Reportar únicamente la cifra más alta sin señalar la divergencia
constituiría una lectura selectiva de la fuente.}

Todas estas curvas señalan el mismo punto metodológico: una única cifra
agregada promedia un régimen en el que el estado dice poco con otro en el que
lo dice casi todo. Ninguno de estos trabajos reporta intervalos de confianza
sobre esas curvas, y ninguno reporta cómo cambia su forma cuando cambia la
población de jugadores.

Hodge \textit{et al.} organizan las características empleadas en la
literatura según el momento en que están disponibles, previas al juego,
durante el juego y posteriores al juego, y recogen la observación de que
las disponibles durante el juego concentran la información más útil para la
predicción \cite{hodge_winprediction_2021}. Su revisión reúne dieciséis
trabajos previos y advierte que la variación en tamaño, composición y versión
del juego entre los conjuntos empleados obliga a tratar con cautela la
comparación de las exactitudes reportadas.

\razon{Esta advertencia debe recogerse de forma explícita porque es el
precedente de dominio de la declaración de comparabilidad que exige la
Sección~\ref{sec:baseline}.}

\subsubsection{Evaluación de estados en juegos de cartas}

Miernik y Kowalski evolucionan funciones lineales y basadas en árboles para
asignar un valor escalar a estados de \emph{Legends of Code and Magic}; su
mejor variante alcanza un 60.2\,\% de tasa de victoria en el torneo de
evaluación, pero la función se optimiza de forma indirecta mediante el
desempeño del agente y no representa una probabilidad aprendida a partir de
pares estado--desenlace \cite{miernik_evolving_2022}.

Xenopoulos \textit{et al.} estiman la probabilidad de victoria por ronda en un
videojuego de disparos tácticos y comparan tres poblaciones de jugadores con
distinto nivel de habilidad \cite{xenopoulos_winprob_2022}. Su representación
del estado combina información global de la ronda con información agregada
por bando, estructura análoga a la que se adopta aquí. Reportan que los
modelos se encuentran bien calibrados dentro de su propio entorno, y que el
entrenado sobre la población de menor nivel presenta mayor pérdida
logarítmica y menor calidad de calibración al evaluarse sobre las poblaciones
profesionales.

\subsubsection{Generalización ante contenido cambiante}

El contenido de estos juegos no se queda quieto. Bertram \textit{et al.} lo
abordan directamente para \emph{Magic: The Gathering}, evaluando mediante la
reserva de un conjunto de cartas completo. Sus resultados separan dos
capacidades que una única puntuación dentro de la distribución confunde: una
codificación por identidad alcanza un 67.80\,\% de acierto sobre datos
reservados del mismo conjunto y no puede producir predicción alguna para
cartas nunca vistas, mientras que las representaciones construidas a partir
de atributos alcanzan un 68.00\,\% dentro de la distribución y un 42.87\,\%
sobre conjuntos no vistos; escalado a un corpus de 2\,990 cartas, el mismo
modelo alcanza un 55.44\,\% sobre un conjunto enteramente reservado
\cite{bertram_cardrepresentations_2024}. La lección trasciende su tarea: una
codificación del contenido basada en la identidad puede ser indistinguible de
una semántica dentro de la distribución y no valer nada fuera de ella.

Chitayat \textit{et al.} proponen una representación de personajes basada en
atributos y habilidades de diseño en lugar de identificadores; un modelo que
la emplea mantiene su desempeño sobre dos versiones posteriores del juego no
utilizadas para entrenar \cite{chitayat_meta_2023}. Su caso de estudio, sin
embargo, predice conteos finales a partir de duración y composición de
equipo, y no probabilidad de victoria desde un estado intermedio.

Angliss \textit{et al.} abordan la diversidad de composiciones. Su banco de
pruebas separa composiciones vistas de composiciones fuera de distribución, y
reporta que el agente entrenado con mayor diversidad transfiere mejor a
composiciones no vistas aun cuando su rendimiento empeora en configuraciones
familiares \cite{angliss_vgcbench_2026}.

Saravanan y Guzdial estudian la evolución del conjunto de estrategias
dominante ante cambios de balance del juego
\cite{saravanan_metadiscovery_2024}. El trabajo se sitúa en el videojuego
competitivo de la franquicia Pok\'emon y no en el juego de cartas
coleccionables, distinción que debe mantenerse presente al interpretar sus
resultados. Los autores definen dicho conjunto como el conocimiento que
excede las reglas del juego y que se manifiesta en los personajes y
estrategias predominantes dentro de la base de jugadores, y observan que
cambia de forma continua.

\razon{La aclaración sobre cuál juego de la franquicia estudia este trabajo
es obligatoria en su primera mención. Presentarlo sin la distinción
atribuiría al artículo un alcance que no tiene.}

\subsubsection{Síntesis del eje}

\noindent\textbf{Enfoques predominantes.} En deportes electrónicos predominan
los métodos supervisados para predicción de resultado y, en trabajos sobre
estado dinámico, las representaciones que incorporan información
intra-partida. En juegos de cartas y en Pok\'emon aparecen además evaluadores
heurísticos, búsqueda, clonación de comportamiento y aprendizaje por
refuerzo, pero dichos enfoques optimizan decisiones o políticas y no
necesariamente producen probabilidades calibradas
\cite{kamal_moba_2025,miernik_evolving_2022,angliss_vgcbench_2026}.

\noindent\textbf{Aspectos considerados resueltos.} La predicción del
desenlace a partir de un estado intermedio es un problema supervisado bien
planteado y tratable en juegos de cartas coleccionables y en juegos de
estrategia en tiempo real; las características tabulares derivadas del estado
transportan señal sustancial
\cite{janusz_hearthstone_2017,bialecki_sc2egset_2023}; y la capacidad
predictiva varía sistemáticamente a lo largo de la partida
\cite{tian_hero_2022}. Existen asimismo protocolos explícitos para estudiar
la robustez ante actualizaciones del juego y ante composiciones no vistas
\cite{chitayat_meta_2023,angliss_vgcbench_2026,bertram_cardrepresentations_2024}.

\noindent\textbf{Aspectos abiertos.} Si las probabilidades estimadas están
calibradas y no meramente bien ordenadas; cómo se comporta el perfil por
etapa bajo un cambio declarado en la población de jugadores o de mazos; y si
algo de esto se sostiene en el \emph{Pok\'emon TCG}, para el que no existe
estudio publicado. Tampoco resulta válido asumir que los resultados obtenidos
en juegos de arena multijugador, en el videojuego competitivo de Pok\'emon o
en \emph{Hearthstone} se transfieren de forma directa, dado que difieren las
reglas, las unidades de observación, los protocolos y la población.

% 
% EJE B
% 
\subsection{Eje B: las técnicas del \textit{track} seleccionado}
\label{subsec:ejeb}

El \textit{track} seleccionado corresponde al aprendizaje automático clásico
sobre datos tabulares. Este apartado establece el estado actual de dichos
métodos, los supuestos bajo los cuales operan, y las condiciones que
determinan la validez de las cifras que producen.

\subsubsection{Estado actual de los métodos tabulares}

Borisov \textit{et al.} presentan una revisión sistemática de los enfoques de
aprendizaje profundo para datos tabulares heterogéneos, junto con una
comparación experimental frente a métodos clásicos bajo un protocolo
uniforme \cite{borisov_tabular_2024}. Los autores organizan la literatura en
tres familias: métodos de transformación de datos, arquitecturas
especializadas y modelos de regularización. Su evaluación reporta que los
conjuntos de árboles con potenciación del gradiente obtienen los mejores
resultados en la mayoría de los conjuntos considerados, y concluye que la
conveniencia de emplear técnicas actuales de aprendizaje profundo sobre datos
tabulares admite en general una respuesta negativa, particularmente para
conjuntos heterogéneos de tamaño reducido.

Grinsztajn \textit{et al.} amplían la evidencia con una comparación sobre
cuarenta y cinco conjuntos de datos, contabilizando el costo de la selección
de hiperparámetros, y reportan que los modelos basados en árboles resultan
superiores para todo presupuesto de búsqueda considerado
\cite{grinsztajn_trees_2022}.

\subsubsection{El rendimiento marginal de la complejidad en este dominio}

La evidencia procedente del propio dominio de juegos apunta en la misma
dirección, y lo hace con cifras. En la competencia de \emph{Hearthstone} se
registraron 296 equipos y 188 enviaron solución; todos los equipos mejor
clasificados usaron redes neuronales, y aun así la puntuación ganadora de
0.8019 supera al resultado de referencia de los organizadores, 0.7846, una red totalmente conectada de dos capas ocultas sobre la representación
tabular, por menos de un dos por ciento, y los organizadores lo hacen
notar explícitamente \cite{janusz_hearthstone_2017}. En SC2EGSet, tres
clasificadores convencionales con hiperparámetros próximos a los
predeterminados quedan a menos de 0.016 de acierto unos de otros: máquina de
vectores de soporte en 0.9071, potenciación del gradiente en 0.8924 y
regresión logística en 0.8916 \cite{bialecki_sc2egset_2023}.

Dos estudios independientes sobre dos juegos distintos coinciden, por tanto,
en que la mayor parte de la señal extraíble de una representación tabular del
estado es alcanzable sin arquitecturas profundas, y en que el residuo es lo
bastante pequeño como para que la interpretabilidad y el costo dominen la
elección.

Karbalaie \textit{et al.} afinan el argumento en dirección a los modelos más
simples, y por una razón inesperada: compararon diez clasificadores y
encontraron que los modelos lineales fueron los más honestos bajo validación
consciente de la agrupación. La regresión logística y el análisis
discriminante lineal difirieron a lo sumo en 0.03 de acierto entre el
protocolo contaminado y el correcto, mientras que los conjuntos de árboles
mostraron brechas de al menos 0.34 y fueron los modelos cuyas puntuaciones
más infló el protocolo defectuoso \cite{karbalaie_participantaware_2026}.

\razon{Este hallazgo debe consignarse porque invierte una expectativa: los
modelos de alta capacidad no son solo más difíciles de interpretar, son los
modelos a los que un protocolo defectuoso favorece. Constituye un argumento
para mantener un modelo lineal regularizado en la comparación como punto de
referencia, y no por formalidad.}

\subsubsection{Supuestos bajo los que operan estos métodos}

Grinsztajn \textit{et al.} identifican tres propiedades de los datos
tabulares que explican la diferencia de desempeño
\cite{grinsztajn_trees_2022}. En primer lugar, las funciones objetivo
asociadas a estos conjuntos no son suaves, mientras que las redes neuronales
presentan un sesgo hacia soluciones suaves; los modelos basados en árboles,
que aprenden funciones constantes por tramos, no exhiben dicho sesgo. En
segundo lugar, los conjuntos tabulares contienen una proporción apreciable de
características no informativas, ante las cuales los modelos neuronales de
tipo perceptrón multicapa resultan menos robustos. En tercer lugar, las
características de un conjunto tabular poseen significado individual,
propiedad que los procedimientos invariantes ante rotaciones no pueden
aprovechar.

\pc{Párrafo argumentando cada una de las tres propiedades sobre el dominio
concreto: la irregularidad de la función objetivo asociada a los umbrales
duros del juego, la presencia previsible de características poco informativas
en el diccionario, y el significado individual de las variables del estado.}

El supuesto que todos estos métodos comparten es que las muestras son
independientes, y es exactamente el supuesto que este corpus viola. Karbalaie
\textit{et al.} cuantifican la consecuencia sobre 623 ensayos de 72
participantes: la validación cruzada estratificada de diez pliegues da un
acierto de 0.91, y el mismo modelo bajo validación dejando un participante
fuera da 0.66 \cite{karbalaie_participantaware_2026}. Su explicación es la que
se transfiere: en datos de medidas repetidas el tamaño muestral efectivo se
aproxima al número de participantes y no al número de ensayos. Sustituir
participante por episodio e instantánea por ensayo enuncia la restricción que
gobierna el presente trabajo. Reportan además que la validación cruzada
agrupada de $k$ pliegues coincide con el diseño de dejar un grupo fuera
dentro de 0.02 de acierto y a un costo mucho menor, lo que resulta decisivo
cuando la unidad de agrupación se cuenta en decenas de miles.

\razon{Brecha de dominio, declarada de forma explícita: se trata de un estudio
de biomecánica, citado por su argumento estadístico sobre observaciones
dependientes y no como resultado del dominio de juegos.}

\noindent\textbf{Límites del alcance de la evidencia tabular.} Borisov
\textit{et al.} reportan que, en el mayor de los conjuntos evaluados, con
características predominantemente continuas, una arquitectura basada en
mecanismos de atención supera a los métodos clásicos
\cite{borisov_tabular_2024}. Esta observación resulta pertinente dado el
volumen de observaciones disponible aquí. Sin embargo, la condición relevante
no es el número de filas sino el número de unidades estadísticamente
independientes, 52\,085 episodios frente a 7\,038\,378 filas, y las
variables del estado son heterogéneas y no predominantemente continuas.

Por su parte, Grinsztajn \textit{et al.} excluyen de su comparación los
conjuntos con datos faltantes, aquellos con variables categóricas de
cardinalidad superior a veinte, los conjuntos cuyas observaciones no son
independientes, y los conjuntos deterministas procedentes de juegos
\cite{grinsztajn_trees_2022}. El problema abordado aquí no es determinista:
el desenlace de una partida no se deduce del estado intermedio, sino que
depende del orden restante del mazo, de la información oculta del oponente y
de las decisiones posteriores de ambos jugadores. Es precisamente esa
condición la que motiva estimar una probabilidad en lugar de deducir una
etiqueta. Respecto del supuesto de independencia, este trabajo adopta el
episodio como unidad independiente.

\razon{Ambos límites deben declararse en el texto y no omitirse. Son las dos
objeciones más directas que puede recibir la justificación del
\textit{track}, y responderlas de forma anticipada es preferible a que surjan
durante la sustentación.}

\subsubsection{Dificultad intrínseca de la tarea}

Brill \textit{et al.} examinan la dificultad de estimar la probabilidad de
victoria mediante un estudio de simulación en el cual la probabilidad
verdadera es calculable de forma exacta \cite{brill_winprob_2026}. Reportan
que la estructura de dependencia propia de los datos observacionales
incrementa el sesgo y la varianza de los estimadores y reduce el tamaño de
muestra efectivo; que resulta preferible conservar la totalidad de las
observaciones por partida antes que retener una sola, dado que el conjunto
completo aporta información sobre la estructura del espacio de covariables; y
que los intervalos de confianza construidos mediante remuestreo presentan
cobertura inferior a la nominal, incluso cuando el remuestreo respeta la
estructura de agrupamiento.

Los autores declaran de forma explícita que sus hallazgos resultan aplicables
a cualquier conjunto de datos en el cual la variable de resultado corresponda
al desenlace de una unidad y las observaciones consistan en las unidades que
conducen a dicho desenlace \cite{brill_winprob_2026}. El corpus empleado aquí
satisface esa descripción.

\subsubsection{La fuga de datos como objeto formal}

Kaufman \textit{et al.} aportan el marco que vuelve verificables, y no
meramente intuitivas, las decisiones de partición. Definen la fuga de datos
como una propiedad intrínseca de la relación entre las entradas
observacionales de un modelo y la instancia objetivo: una característica es
legítima solo si es observable para el propósito de inferir ese objetivo en
particular, evaluada instancia por instancia
\cite{kaufman_leakage_2011}. Dos consecuencias gobiernan el diseño de este
trabajo. La primera es el requisito de la no máquina del tiempo: un modelo
legítimo solo puede usar información de un instante no posterior al del
objetivo. La segunda es la separación aprender--predecir, según la cual cada
observación se etiqueta con la información necesaria para decidir su
legitimidad y el conjunto de entrenamiento se construye después de modo que
solo entren entradas legítimas.

Advierten además que la fuga entra por las decisiones de diseño, diseño de
características, selección de modelo, elección de hiperparámetros informada
por conocimiento que el modelo desplegado no tendría, a la que llaman la
forma más peligrosa porque un evaluador no puede detectarla ni siquiera
buscándola, y señalan que eliminar las fugas evidentes puede no sellar el
resto.

Sasse \textit{et al.} actualizan ese marco convirtiéndolo en un catálogo de
escenarios organizados como fuga de prueba a entrenamiento, de prueba a
prueba, de característica a objetivo, de objetivo, de conjunto de datos y de
variables de confusión \cite{sasse_leakage_2025}. Tres de sus entradas se
corresponden con decisiones de este trabajo. Su tratamiento de la fuga de
prueba a entrenamiento enumera el no tener en cuenta muestras dependientes
como causa primaria. Su caso de preprocesamiento establece que el escalado,
la imputación, la reducción de dimensionalidad y el sobremuestreo estimados
sobre el conjunto completo antes de particionar constituyen por sí mismos
fuga de datos. Y su categoría de característica a objetivo cubre el caso que
aquí más importa: variables informativas y disponibles durante el aprendizaje
pero ausentes en el despliegue, cuya prescripción es una pregunta antes que
una regla, comprobar si cada característica estaría disponible tras el
despliegue.

Figueiredo y Mendes complementan esa formalización con una medición. Estudian
el efecto de la partición aleatoria sobre conjuntos de imágenes con alta
correlación espaciotemporal, y reportan que produce una inflación apreciable
de las métricas respecto de una partición que asigna agrupaciones completas
\cite{figueiredo_leakage_2024}. Su trabajo pertenece al dominio de la visión
por computador y se considera aquí un análogo metodológico; el mecanismo
subyacente, observaciones casi idénticas a ambos lados de la partición,
es independiente del dominio. Cabe destacar que su contribución central
consiste en un procedimiento para inferir las agrupaciones, dado que las
imágenes no vienen etiquetadas por escena; en este corpus la agrupación viene
dada de forma exacta por el identificador de episodio.

\subsubsection{Evaluación de la calidad de las probabilidades}

Silva Filho \textit{et al.} presentan una revisión de los métodos para
evaluar y mejorar las probabilidades predichas por un clasificador
\cite{silvafilho_calibration_2023}. Un modelo se considera calibrado cuando
las probabilidades que predice se corresponden con las frecuencias
observadas. Los autores exponen que las reglas de puntuación propias se
descomponen en un componente de calibración y uno de refinamiento,
descomposición que fundamenta el conjunto de métricas adoptado aquí. Señalan
asimismo que los estimadores del error de calibración basados en
discretización dependen del número y la construcción de los intervalos
empleados.

Niculescu-Mizil y Caruana examinan las probabilidades producidas por distintas
familias de algoritmos y reportan que los métodos de máximo margen, entre
ellos los conjuntos de árboles con potenciación, desplazan la masa de
probabilidad alejándola de los extremos, produciendo una distorsión de forma
sigmoidea \cite{niculescumizil_probabilities_2005}. Reportan igualmente que
la regresión logística y los conjuntos por agregación resultan bien
calibrados sin corrección posterior, y que recalibrar estos últimos puede
deteriorarlos. Su análisis de curvas de aprendizaje establece que el método
de recalibración apropiado depende del tamaño del conjunto disponible para
dicho ajuste.

\razon{Esta referencia debe aparecer porque convierte la calibración de una
preocupación general en una predicción concreta sobre los modelos que este
trabajo va a emplear. Sin ella, la sección de métricas carecería de
justificación específica.}

Ojeda \textit{et al.} comparan de forma sistemática los enfoques de
recalibración disponibles para la estimación de probabilidades de dos clases,
evaluando entre otros aspectos su capacidad de generalizar estimaciones
precisas a poblaciones externas \cite{ojeda_calibrating_2023}. Dimitriadis
\textit{et al.} abordan la inestabilidad de los diagramas de fiabilidad
construidos mediante discretización, y proponen un procedimiento que genera
diagramas reproducibles junto con una medida numérica de descalibración
\cite{dimitriadis_reliability_2021}.

\decidir{Las recomendaciones de Ojeda \textit{et al.} y las de
Niculescu-Mizil y Caruana no coinciden de forma inmediata respecto del método
de recalibración preferible. La resolución requiere revisar si el escenario
evaluado por Ojeda \textit{et al.}, transferencia a poblaciones
externas, corresponde al protocolo temporal de este trabajo, en cuyo caso
su recomendación tendría precedencia. La decisión debe tomarse antes de la
Etapa 2.}

\subsubsection{Requisitos metodológicos del \textit{track}}

Emmanuel \textit{et al.} revisan los mecanismos de datos faltantes y muestran
que la eficacia de un método de imputación depende de las características del
conjunto: en su experimento, la superioridad relativa entre procedimientos se
invierte según el conjunto evaluado \cite{emmanuel_missing_2021}. El
resultado relevante no es que una técnica resulte superior, sino que la
estrategia depende del mecanismo de ausencia. Una frecuencia elevada de
valores ausentes no justifica por sí sola imputarlos de forma automática.

Van den Goorbergh \textit{et al.} estudian modelos de regresión logística bajo
ausencia de corrección, submuestreo aleatorio, sobremuestreo aleatorio y
generación sintética de la clase minoritaria. Las correcciones no mejoran de
forma sistemática el área bajo la curva, y sí producen una sobreestimación
pronunciada de la probabilidad de la clase minoritaria, con el consiguiente
deterioro de la calibración \cite{vandenGoorbergh_imbalance_2022}.

\razon{El alcance de esta última referencia debe respetarse. El trabajo no
evalúa la ponderación de clases ni familias distintas de la regresión
logística, de modo que no justifica preferir la ponderación como regla
general. Presentarlo de otro modo excedería lo que la fuente sostiene.}

\subsubsection{Justificación técnica del \textit{track}}

El Track A se elige por la fuerza de esa evidencia y no por descarte. El
corpus de repeticiones se reduce de forma natural a una tabla de ancho fijo
con campos numéricos, binarios y categóricos, que es el régimen en el que los
estudios del dominio de juegos citados muestran una paridad casi total entre
modelos clásicos y profundos
\cite{janusz_hearthstone_2017,bialecki_sc2egset_2023}, y en el que la
evidencia tabular general favorece a los conjuntos de árboles
\cite{borisov_tabular_2024,grinsztajn_trees_2022}.

La pregunta científica es cuáles características del estado transportan
información sobre la victoria, lo que exige modelos cuyas atribuciones puedan
inspeccionarse. Y el riesgo dominante es de protocolo y no de capacidad: la
contaminación medida de una partición a nivel de fila hace que sea el diseño
de la partición, y no la arquitectura, la variable que determina si un
resultado tiene significado.

Un modelo secuencial profundo sobre la trayectoria, la alternativa natural
del Track B, respondería una pregunta distinta y más amplia, heredando todos
los riesgos de fuga y ninguna de las ventajas de interpretabilidad. El
aprendizaje por refuerzo, lectura natural de un corpus de episodios, se
descarta porque el objeto de estudio es la función de valor de forma aislada,
aprendida \emph{offline} a partir de desenlaces registrados, y no una
política.

%
% EJE C
% 
\subsection{Eje C: el conjunto de datos y sus análogos}
\label{subsec:ejec}

\subsubsection{Ausencia de literatura previa sobre este corpus}

No localizamos ninguna publicación arbitrada que utilice los episodios del
PTCG AI Battle Challenge \cite{ptcg_dataset_2026}. El corpus se publicó en
2026 y está ligado a una competencia todavía en curso. De acuerdo con el
enunciado del curso, el eje se cubre por tanto con conjuntos de datos
análogos del mismo dominio, y la sustitución se argumenta en lugar de darse
por supuesta.

\razon{La ausencia debe declararse de forma explícita en este apartado y
recogerse después como brecha. Declararla convierte una carencia en un
hallazgo de la revisión.}

\subsubsection{Por qué los análogos son comparables, y dónde no lo son}

SC2EGSet es el análogo estructural más cercano: un corpus de repeticiones
competitivas de un juego de dos jugadores, procesado en registros de estado
por partida, publicado junto con sus herramientas de extracción y acompañado
de experimentos de predicción de victoria sobre características tabulares
\cite{bialecki_sc2egset_2023}. Comparte la forma de la tubería, repetición cruda, estado analizado, características tabulares, etiqueta
supervisada de desenlace, y el hecho de que los datos fueron producidos por
competencia y no por experimento. Difiere en género, en tanto
\emph{StarCraft~II} es un juego de estrategia en tiempo real y no un juego de
cartas por turnos con recursos discretos, y difiere en unidad de observación,
puesto que sus modelos consumen agregados por partida y no estados
instantáneos.

El conjunto de la competencia de \emph{Hearthstone} es el análogo más cercano
en género: un juego de cartas coleccionables con mazos, mano, tablero e
información oculta, cuya tarea es precisamente la predicción del ganador a
partir de un estado intermedio \cite{janusz_hearthstone_2017}. Difiere en la
calidad de los agentes que lo generaron y en que sus estados fueron
publicados por los organizadores en formato tabular en lugar de ser
analizados por quien modela.

El corpus de selección de cartas de \emph{Magic: The Gathering} es análogo
solo en el sentido estrecho que importa para una de nuestras decisiones: es
un juego de cartas coleccionables cuyo contenido cambia con el tiempo y cuya
evaluación reserva contenido no visto
\cite{bertram_cardrepresentations_2024}. Su tarea es predecir una elección
humana de carta, no un desenlace de partida, y no debe leerse como un
resultado de predicción de victoria.

\subsubsection{Problemas documentados de los corpus comparables}

Los problemas que estos autores declaran son los problemas que deberíamos
esperar.

Bia\l ecki \textit{et al.} afirman que su corpus contiene únicamente
repeticiones de nivel de torneo y que, en consecuencia, no admite análisis
entre niveles de habilidad; que los cambios del juego a lo largo del tiempo
hacen que algunos campos tomen valor cero, de modo que la representación
cambió con el tiempo; que la completitud dependió de qué organizadores
decidieron publicar; y que deliberadamente no filtraron repeticiones
terminadas de forma prematura, con el fin de preservar las distribuciones
originales \cite{bialecki_sc2egset_2023}.

Cada uno de esos puntos tiene su contraparte medida en nuestro corpus: los
episodios se seleccionan por puntuación media de los agentes y se truncan a
un tope diario de volumen; seis versiones del motor coexisten en la muestra
auditada, de la 1.30.1 a la 1.32.4; y los episodios abandonados están
presentes y son identificables.

Janusz \textit{et al.} documentan dos problemas propios. Sus estados provenían
de partidas entre agentes que elegían sus decisiones al azar, lo que
constituye una población específica y no representativa; y se detectó una
fuga de datos involuntaria durante las primeras semanas de la competencia,
que obligó a los organizadores a mover 1\,250\,000 casos del conjunto de
prueba al de entrenamiento \cite{janusz_hearthstone_2017}.

\razon{Vale la pena enunciar lo segundo con claridad: un incidente de fuga
dentro de una competencia organizada profesionalmente sobre exactamente esta
tarea es la evidencia más fuerte disponible de que el riesgo es real y no
hipotético.}

\subsubsection{Cambio de distribución y qué se degrada primero}

La respuesta parcial de Janusz \textit{et al.} al problema de la población es
también el resultado de generalización más útil del que disponemos. Habiendo
entrenado sobre estados de juego aleatorio, construyeron un segundo corpus de
28\,604 partidas y 814\,407 estados generados por agentes fuertes de búsqueda
en árbol Monte Carlo, y observaron que el área bajo la curva total caía solo
de aproximadamente 0.79 a aproximadamente 0.75 al evaluar sobre esa población
desplazada \cite{janusz_hearthstone_2017}. Un cambio en la calidad de los
agentes cuesta, por tanto, algo, pero mucho menos de lo que podría temerse.

Lo que ese experimento no mide es la calibración, y la evidencia sistemática
indica que esa es la magnitud en mayor riesgo. Guo \textit{et al.} examinaron
4\,457 referencias y retuvieron quince estudios que cuantificaban el cambio
temporal e implementaban alguna mitigación. Los once estudios que evaluaron
calibración reportaron deterioro de la calibración, mientras que solo tres de
los doce que evaluaron discriminación reportaron deterioro de la
discriminación \cite{guo_temporalshift_2021}. Su inventario de mitigaciones
está dominado por el reajuste del modelo, la calibración de probabilidades,
la actualización del modelo y la selección de modelo, frente a solo dos
enfoques a nivel de características, y su síntesis es que todas esas
estrategias preservaron la calibración pero ninguna preservó uniformemente la
discriminación.

Aportan además el vocabulario que este trabajo necesita, distinguiendo el
cambio de covariables ---$P_{\text{ent}}(x)\neq P_{\text{prueba}}(x)$ con
$P_{\text{ent}}(y|x)=P_{\text{prueba}}(y|x)$--- del cambio de concepto
---$P_{\text{ent}}(y|x)\neq P_{\text{prueba}}(y|x)$ con
$P_{\text{ent}}(x)=P_{\text{prueba}}(x)$---, y señalando que ambos pueden
coexistir.

\razon{Brecha de dominio, declarada de forma explícita: se trata de
predicción clínica, citada por la asimetría que mide y por la terminología
que estandariza, no como hallazgo del dominio de juegos. Sus propias
limitaciones declaradas, barrido incompleto de actas de congresos y
ausencia de tratamiento del cambio geográfico, no afectan a ninguno de los
dos usos.}

% 
% TABLA COMPARATIVA
% 
\subsection{Análisis comparativo}
\label{subsec:comparativo}

\safeinput{tab/tabla_comparativa}

De la Tabla~\ref{tab:comparativa} emergen cuatro patrones.

\noindent\textbf{Una división de métricas que bloquea la comparación
directa.} Los estudios sobre estados de juego reportan escalares distintos, área bajo la curva en unos casos, exactitud en otros, y ninguno reporta
una regla de puntuación propia ni una curva de fiabilidad, con la excepción
de Xenopoulos \textit{et al.}, que sí incorporan una medida de error de
calibración \cite{xenopoulos_winprob_2022}. Cualquier afirmación de que uno
de estos sistemas está mejor calibrado que otro carecería de respaldo, porque
la mayoría no lo midió.

\noindent\textbf{La generalización se mide contra cosas diferentes.} Calidad
de los agentes \cite{janusz_hearthstone_2017}, conjuntos de cartas no vistos
\cite{bertram_cardrepresentations_2024}, identidad del participante
\cite{karbalaie_participantaware_2026}, tiempo calendario
\cite{guo_temporalshift_2021}, versiones del juego \cite{chitayat_meta_2023},
composiciones de equipo \cite{angliss_vgcbench_2026}, nivel de habilidad
\cite{xenopoulos_winprob_2022}. Y las degradaciones reportadas son en todos
los casos lo bastante grandes como para importar: 0.79 a 0.75 de área bajo la
curva, 67.20\,\% a 55.44\,\% de acierto, 0.91 a 0.66 de acierto. Una única
cifra dentro de la distribución es la magnitud menos informativa que reporta
cualquiera de estos trabajos. Un único holdout aleatorio sería, por tanto,
insuficiente para sostener una conclusión firme de generalización.

\noindent\textbf{La unidad de observación es menor que la unidad
independiente, y los trabajos difieren en si actúan en consecuencia.} Janusz
\textit{et al.} separan entrenamiento y prueba a nivel de partidas completas,
de modo que su protocolo está agrupado aunque no lo formulen así. Karbalaie
\textit{et al.} convierten la agrupación en el objeto de estudio. Bia\l ecki
\textit{et al.} reportan validación cruzada de cinco pliegues sobre 22\,230
muestras promediadas derivadas de 11\,115 partidas únicas, dos filas por
partida, una por jugador, sin describir restricción alguna de agrupación a
nivel de partida.

\razon{Lo siguiente es una inferencia nuestra y no una afirmación de los
autores: si esos pliegues se extrajeron a nivel de fila, ambas filas de una
partida podrían caer a lados opuestos de una partición, y sus aciertos
reportados serían optimistas en una cantidad desconocida. Se señala como
salvedad de comparabilidad y no como error, puesto que el artículo no declara
el particionador que utilizó.}

\noindent\textbf{El resultado depende de una población acotada.} Los estudios
sobre juegos de arena multijugador revisados se restringen a jugadores de
nivel elevado o profesional; el trabajo sobre el videojuego competitivo de
Pok\'emon opera con un entorno y un formato distintos de los del juego de
cartas; y los dos corpus de repeticiones declaran esa limitación para sí
mismos. Ello obliga a declarar los límites de extrapolación y a evitar
presentar un resultado obtenido sobre agentes artificiales como evidencia
acerca de jugadores humanos.

% BRECHAS
\subsection{Brechas identificadas}
\label{subsec:brechas}

\noindent\textbf{Carencias encontradas en la literatura.}

\begin{enumerate}
  \item \textbf{Ningún trabajo arbitrado utiliza este corpus.} Los episodios
        del PTCG AI Battle Challenge no tienen descripción publicada, ni
        diccionario de datos documentado, ni resultado de referencia
        reportado \cite{ptcg_dataset_2026}.
  \item \textbf{La tarea no se ha estudiado en el \emph{Pok\'emon TCG}.} La
        predicción de victoria desde el estado está establecida en
        \emph{Hearthstone} \cite{janusz_hearthstone_2017} y en
        \emph{StarCraft~II} \cite{bialecki_sc2egset_2023}, pero no en un
        juego cuya estructura de recursos se apoya en cartas de premio,
        asignación de energía, banca y evolución.
  \item \textbf{Se reporta discriminación sin calibración.} La literatura de
        predicción de desenlace revisada reporta predominantemente métricas
        de ordenamiento o exactitud
        \cite{janusz_hearthstone_2017,bialecki_sc2egset_2023,tian_hero_2022,hodge_winprediction_2021},
        con la excepción de Xenopoulos \textit{et al.}, que incorporan una
        medida de error de calibración \cite{xenopoulos_winprob_2022}.
        Karbalaie \textit{et al.} colocan explícitamente la calibración fuera
        de su alcance y la recomiendan como trabajo futuro
        \cite{karbalaie_participantaware_2026}, y Guo \textit{et al.}
        muestran que es la magnitud que se degrada primero
        \cite{guo_temporalshift_2021}. Un modelo pensado como evaluador de
        posiciones está siendo validado, por tanto, sobre la propiedad
        equivocada en la mayoría de los casos.
  \item \textbf{Evaluación de métricas probabilísticas pendiente en la
        comparación de modelos tabulares.} Grinsztajn \textit{et al.}
        identifican explícitamente esta línea como no cubierta por su
        comparación \cite{grinsztajn_trees_2022}.
  \item \textbf{La evaluación agrupada se aplica de forma inconsistente en
        corpus de estados de juego.} El costo medido de ignorar la
        dependencia solo está disponible desde otros campos
        \cite{karbalaie_participantaware_2026,sasse_leakage_2025,figueiredo_leakage_2024};
        dentro de los conjuntos de datos de juegos la restricción unas veces
        se satisface sin ser nombrada \cite{janusz_hearthstone_2017} y otras
        queda sin describir \cite{bialecki_sc2egset_2023}.
  \item \textbf{La generalización rara vez se evalúa bajo un cambio
        declarado, y nunca bajo los tres a la vez.} Cuando se hace, la
        degradación es sustancial
        \cite{janusz_hearthstone_2017,bertram_cardrepresentations_2024}, y la
        evidencia sistemática encuentra solo quince estudios que a la vez
        cuantifiquen el cambio temporal y actúen sobre él
        \cite{guo_temporalshift_2021}. No existe un protocolo único que
        combine partida nueva, composición no vista y periodo posterior para
        esta tarea.
\end{enumerate}

\razon{La brecha tercera no puede formularse como ausencia absoluta. Al menos
un trabajo de la propia bibliografía reporta calibración, de modo que afirmar
que ningún trabajo lo hace sería contradicho por una referencia citada en
este mismo documento.}

\noindent\textbf{Brechas que este trabajo aborda.} Las brechas 2, 3, 5 y 6.
La segunda, planteando la tarea en el \emph{Pok\'emon TCG} por primera vez y
reportando el perfil por turno de la capacidad predictiva. La tercera,
reportando pérdida logarítmica, puntaje de Brier y curvas de fiabilidad junto
al área bajo la curva, y poniendo a prueba la hipótesis H3 de forma
explícita. La quinta, agrupando por identificador de episodio toda partición,
todo intervalo de confianza y todo pliegue de validación cruzada, y
reportando la contaminación medida que produciría la alternativa ingenua. La
sexta, evaluando bajo tres cambios declarados, episodios no vistos, mazos
no vistos y periodos posteriores, y reportando cada uno por separado en
lugar de promediarlos.

\noindent\textbf{Brechas que quedan fuera del alcance.} La primera se
documenta pero no se cierra: un proyecto de un semestre no puede establecer
una referencia arbitrada para el corpus, y este trabajo no pretende serlo. La
cuarta excede el alcance, dado que este no es un estudio comparativo de
familias de modelos sobre múltiples conjuntos de datos. La comparabilidad
estricta con el antecedente de \emph{Hearthstone} queda igualmente fuera de
alcance por las razones enumeradas en la Sección~\ref{sec:baseline}: las
diferencias son estructurales y no pueden eliminarse reejecutando ninguno de
los dos experimentos. Más allá de estas, no se persigue construir un agente
competitivo, aprender una política, separar la deriva del conjunto de
estrategias dominante del cambio de versión del motor, la muestra auditada
confunde ambas causas, ni generalizar ningún hallazgo a jugadores humanos.

\razon{Reconocer explícitamente las brechas que el trabajo no aborda es un
requisito del enunciado y no una debilidad. Un apartado que declarara
abordarlas todas resultaría menos creíble que uno que delimita su alcance.}